% overleaf-overcharge.tex — Prueba de esfuerzo para el compilador
% Eric Lucero González — https://ericlucero.dev — Licencia CC BY-NC-SA 4.0
%
% Dibuja el conjunto de Mandelbrot píxel a píxel con TikZ. Cada píxel repite
% z = z^2 + c hasta \iteraciones veces, y TeX hace esas cuentas con macros:
% es lentísimo a propósito. El tiempo crece con el cuadrado de \resolucion.
%
% Cómo usarlo:
%   1. Compila con pdfLaTeX (en Overleaf es el compilador por defecto).
%   2. Al final del PDF aparece cuánto ha tardado la compilación.
%   3. Cambia \resolucion y mira dónde se rinde cada plan de Overleaf.
%   El tiempo crece con el cuadrado: el doble de resolución, el cuádruple de tiempo.
\documentclass[a4paper]{article}
\usepackage[T1]{fontenc}
\usepackage[spanish]{babel}
\usepackage{lmodern}
\usepackage[margin=2cm]{geometry}
\usepackage{amsmath, amssymb}
\usepackage{tikz}
\pagestyle{empty}

% ── Carga: cambia solo estas dos líneas ──────────────────────────
\def\resolucion{100}   % píxeles por lado (60 = 3.600 píxeles)
\def\iteraciones{30}  % iteraciones máximas por píxel

% ── Región del plano complejo que se dibuja ──────────────────────
\def\xmin{-2.2} \def\xmax{0.8}
\def\ymin{-1.5} \def\ymax{1.5}
\def\lado{14cm}       % tamaño del dibujo

% ── Iteración de un píxel ────────────────────────────────────────
% Deja en \iter cuántas iteraciones hacen falta para que |z| > 2
% (o \iteraciones si el punto parece estar dentro del conjunto).
\newcount\iter
\newif\ifsigue
\newcommand\mandelbrot[2]{% #1 = parte real de c, #2 = parte imaginaria
  \def\zr{0}\def\zi{0}\iter=0 \siguetrue
  \loop
  \pgfmathsetmacro\zrr{\zr*\zr}%
  \pgfmathsetmacro\zii{\zi*\zi}%
  \pgfmathsetmacro\modulo{\zrr+\zii}%
  \ifdim\modulo pt>4pt \siguefalse \else
  \ifnum\iter<\iteraciones\relax
  \pgfmathsetmacro\zi{2*\zr*\zi+#2}%
  \pgfmathsetmacro\zr{\zrr-\zii+#1}%
  \advance\iter by 1
  \else \siguefalse \fi
  \fi
  \ifsigue \repeat
}

\begin{document}

\section*{Prueba de esfuerzo: el conjunto de Mandelbrot hecho con \TeX}

Este dibujo tiene \(\resolucion \times \resolucion\) píxeles y hasta
\iteraciones{} iteraciones por píxel. Cada cuadrado se calcula con macros de
\TeX{} y Ti\textit{k}Z, sin programas externos: el compilador hace todas las
cuentas. Cambia \verb|\resolucion| en el preámbulo para aumentar la carga.

\begin{center}
  \begin{tikzpicture}
    % ── Preparación ───────────────────────────────────────────────
    % La rejilla tiene \resolucion × \resolucion píxeles cuadrados.
    % \paso: lado de cada píxel en el papel (\lado / \resolucion, en pt).
    \pgfmathsetmacro\paso{\lado/\resolucion}
    % \dx, \dy: cuánto avanza c (en el plano complejo) de un píxel al siguiente.
    \pgfmathsetmacro\dx{(\xmax-\xmin)/\resolucion}
    \pgfmathsetmacro\dy{(\ymax-\ymin)/\resolucion}
    % Índice del último píxel: los bucles van de 0 a \resolucion - 1.
    \pgfmathtruncatemacro\ultimo{\resolucion-1}

    % ── Recorrido de la rejilla: \i = columna (eje real), \j = fila (eje imaginario)
    \foreach \i in {0,...,\ultimo} {
      \foreach \j in {0,...,\ultimo} {
        % c = \cr + i·\ci, tomado en el centro del píxel (de ahí el +0.5).
        \pgfmathsetmacro\cr{\xmin+(\i+0.5)*\dx}
        \pgfmathsetmacro\ci{\ymin+(\j+0.5)*\dy}
        % Itera z_{n+1} = z_n^2 + c y deja en \iter cuántas vueltas aguantó.
        \mandelbrot{\cr}{\ci}
        % ── Color del píxel ───────────────────────────────────────
        \ifnum\iter=\iteraciones\relax
        % No escapó en \iteraciones vueltas: lo damos por dentro de M.
        \def\relleno{black}%
        \else
        % Escapó en la vuelta \iter: el tono va de 0 (blanco, escapa al
        % instante) a 100 (cian puro, casi no escapa).
        \pgfmathtruncatemacro\tono{100*\iter/\iteraciones}
        \edef\relleno{cyan!\tono!white}%
        \fi
        % ── Dibujo del píxel ──────────────────────────────────────
        % Cuadrado con esquina inferior izquierda en (\i·\paso, \j·\paso).
        % Los +0.1pt solapan un poco los píxeles para que no se vean
        % líneas blancas finas entre ellos en algunos visores de PDF.
        \fill[\relleno] (\i*\paso pt, \j*\paso pt) rectangle ++(\paso pt+0.1pt, \paso pt+0.1pt);
      }
    }
  \end{tikzpicture}
\end{center}

\ifdefined\pdfelapsedtime
\bigskip
\noindent\textbf{Tiempo de compilación hasta aquí:}
\fpeval{round(\pdfelapsedtime/65536, 1)}~segundos
(pdf\LaTeX{}, \(\resolucion \times \resolucion\) píxeles, \iteraciones{} iteraciones).
\fi

% ── La idea matemática ──────────────────────────────────────────
\paragraph{Qué se dibuja.} Para cada punto \(c = x + iy\) del plano complejo se
repite, empezando en \(z_0 = 0\),
\[
  z_{n+1} = z_n^2 + c .
\]
Si la sucesión se mantiene acotada, \(c\) pertenece al conjunto de Mandelbrot
\(\mathcal{M}\); en cuanto \(\lvert z_n\rvert > 2\), se sabe que escapará al
infinito. Separando partes real e imaginaria, \(z = a + bi\):
\[
  a_{n+1} = a_n^2 - b_n^2 + x, \qquad
  b_{n+1} = 2\,a_n b_n + y, \qquad
  \text{escapa si } a_n^2 + b_n^2 > 4 .
\]
Los puntos que no escapan en \iteraciones{} iteraciones se pintan de negro; el
resto, más claros cuanto antes escapan.

\end{document}
